\documentclass[11pt]{article}
\usepackage{amsmath,amssymb,amsthm}

\newtheorem{theorem}{Theorem}
\newtheorem{corollary}[theorem]{Corollary}
\newtheorem{lemma}[theorem]{Lemma}

\newcommand{\Up}{U_+}
\newcommand{\mass}{\operatorname{B}}
\newcommand{\car}{\operatorname{C}}

\title{A phase strip and a Jacobsthal positive-defect suffix bound}
\author{}
\date{27 August 2026}

\begin{document}
\raggedbottom
\maketitle

For a positive odd integer \(x\), put
\[
 \Up(x):=\frac{3x+1}{2^{v_2(3x+1)}}.
 \tag{1}
\]
For a positive exponent word \(v=(a_1,\ldots,a_\ell)\), define
\[
 \mass(v):=\sum_{i=1}^{\ell}a_i,
 \qquad
 \car(v):=\sum_{i=1}^{\ell}
 3^{\ell-i}2^{a_1+\cdots+a_{i-1}}.
 \tag{2}
\]

We compare two exact paths of the same raw length and with the same
endpoint.  At a common boundary, write their state relation as
\[
 z=2^\Delta x+c,
 \qquad
 u:=\frac{c}{2^\Delta},
 \tag{3}
\]
where \(\Delta\) is the difference between the masses of the two
remaining suffixes.  At every record-atom boundary, the established
distortion theorem gives
\[
 \Delta\geq-1,
 \qquad
 -1\leq u\leq5.
 \tag{4}
\]
The parity lattice gives even \(c\) for \(\Delta=0\), and odd \(c\)
for \(\Delta\geq1\).  At \(\Delta=-1\), a nonempty remaining suffix
has \(c=-1/2\).

\section{The phase strip}

Every nontrivial record atom begins with exponent one:
\[
 B=1\,|\,2,\qquad
 C=1\,|\,113,\qquad
 D=1\,|\,122.
 \tag{5}
\]
The vertical bar marks the phase immediately after this first exponent.

Suppose that the right state after one of the tails
\[
 R\in\{2,113,122\}
 \tag{6}
\]
is \((d,c)\), and let \(v\) be the aligned alternative block.  The
state \((\Delta',c')\) before the tail satisfies
\[
 \Delta'=d+\mass(v)-\mass(R),
 \tag{7}
\]
\[
 c'=\frac{2^{\Delta'}\car(R)+2^{\mass(v)}c-\car(v)}
 {3^{|R|}}.
 \tag{8}
\]

\begin{lemma}[lower strip at the first phase]
At the phase marked in (5), every genuine distortion state satisfies
\[
 \boxed{\frac{c'}{2^{\Delta'}}\geq-1.}
 \tag{9}
\]
\end{lemma}

\begin{proof}
First take \(R=2\), and write \(v=(a)\).  If \(d\geq1\), the parity
lattice sharpens (4) to
\[
 u\geq-1+2^{-d}.
 \tag{10}
\]
Since \(\Delta'=d+a-2\), equations (7) and (8) give
\[
 u'=\frac43u+\frac13-\frac{2^{2-d-a}}3
 \geq
 -1+\frac{2^{-d}(4-2^{2-a})}{3}
 \geq-1.
 \tag{11}
\]
For \(d=0\), we have \(u\geq0\), and the same formula gives
\(u'\geq-1/3\).  For \(d=-1\), substitution of \(c=-1/2\) gives
\[
 c'=-2^{a-3}-\frac13,
 \tag{12}
\]
which does not lie in the required dyadic parity lattice.  Hence no
such transition exists.

Now take \(R=113\) or \(R=122\).  Their masses are five and their
carries are respectively \(19\) and \(23\).  Every positive
three-letter word satisfies
\[
 \frac{\car(v)}{2^{\mass(v)}}\leq\frac{19}{8}.
 \tag{13}
\]
For \(d\geq2\), (10) and (8) therefore give
\[
 u'\geq
 \begin{cases}
 (-13-44\cdot2^{-d})/27,&R=113,\\
 (-9-44\cdot2^{-d})/27,&R=122.
 \end{cases}
 \tag{14}
\]
Both lower bounds are at least \(-1\).

It remains to treat \(d=-1,0,1\).  The strip and parity lattice leave
ten right states.  Exact substitution gives the following minimum
values of \(u'\); a dash means that no transition exists:
\[
\begin{array}{c|c|c|c}
d&c&R=113&R=122\\ \hline
-1&-1/2&-7/8&\text{--}\\
0&0&-1&0\\
0&2&2&\text{--}\\
0&4&\text{--}&21/4\\
1&-1&\text{--}&1/16\\
1&1&9/8&\text{--}\\
1&3&19/8&81/32\\
1&5&\text{--}&241/64\\
1&7&4&\text{--}\\
1&9&11/2&11/2.
\end{array}
\tag{15}
\]
For completeness, this table has a finite exact reduction.  When
\(\Delta'<3\), the total mass of \(v\) is fixed.  When
\(\Delta'\geq3\), write \(v=(a,b_1,b_2)\).  Reduce \(b_1,b_2\) to
positive residues modulo
\[
 \operatorname{ord}_{27}(2)=18,
 \tag{16}
\]
and, for each parity, reduce \(a\) to the first positive value with
\(\Delta'\geq3\).  These operations preserve transition
admissibility and can only increase
\(\car(v)/2^{\mass(v)}\), hence can only decrease \(u'\).
The resulting finite table is (15), and every entry is at least
\(-1\).
\end{proof}

\section{The positive-defect suffix}

Define the Jacobsthal number
\[
 J_\delta:=\frac{2^\delta-(-1)^\delta}{3}.
 \tag{17}
\]

\begin{theorem}[Jacobsthal lower bound for a positive-defect suffix]
Let \(W=PU\), with \(P\in\{B,C,D\}\), be an exact record-atom word
from \(q\) to \(y\).  Let \(F\) be an exact word of the same raw length
from \(r\) to \(y\), and suppose
\[
 \mass(F)-\mass(W)=\delta\geq1.
 \tag{18}
\]
Then
\[
 \boxed{
 r\geq2^\delta q-J_\delta
 =
 2^\delta q-\frac{2^\delta-(-1)^\delta}{3}.}
 \tag{19}
\]
\end{theorem}

\begin{proof}
Write
\[
 r=2^\delta q+c.
 \tag{20}
\]
Both \(q\) and \(r\) are odd, so \(c\) is odd.  Let \(a\) be the first
exponent of \(F\).  Since the first exponent of \(W\) is one, the
distortion state immediately after these two steps is
\[
 d=\delta+1-a,
 \qquad
 c_1=\frac{3c+1-2^\delta}{2^a}.
 \tag{21}
\]
Lemma~1 gives \(c_1/2^d\geq-1\).  Therefore
\[
 3c+1-2^\delta\geq-2^{a+d}=-2^{\delta+1},
 \tag{22}
\]
and hence
\[
 3c\geq-2^\delta-1.
 \tag{23}
\]
If \(\delta\) is odd, (23) gives \(c\geq-J_\delta\) directly.  If
\(\delta\) is even, it gives
\[
 c\geq-J_\delta-\frac23.
 \]
Since \(c\) and \(J_\delta\) are odd integers, this again implies
\(c\geq-J_\delta\).  Substitution in (20) proves (19).
\end{proof}

\begin{corollary}[improved fixed-seed depth threshold]
Let \(n\) be the least positive odd integer with an unbounded
accelerated orbit.  If a variable-depth competitor uses a suffix of
positive defect \(\delta\), put
\[
 r_\delta:=2^\delta q-J_\delta.
 \tag{24}
\]
If its inverse prefix has length \(h\), excess \(t\geq T\), and start
below \(n\), then necessarily
\[
 \boxed{
 h>
 \log_{3/2}\left(
 \frac{r_\delta+1}
 {\displaystyle 2/3+(n+1/3)/2^T}
 \right).}
 \tag{25}
\]
For \(T=3\), the denominator in (25) is
\((3n+17)/24\).
\end{corollary}

\begin{proof}
Theorem~2 supplies the suffix lower bound \(r\geq r_\delta\).  Apply
the sharp all-excess prefix interval with \(r_\delta\) in place of the
suffix start and solve for \(h\).
\end{proof}

The estimate (19) replaces the earlier bound \(r\geq2^\delta(q-1)\).
It removes approximately two thirds of the former uncertainty in the
negative offset.  The deeper positive-defect prefix branch is more
constrained, but it is not yet excluded.

\end{document}
